\chapter*{Conclusion}\addcontentsline{toc}{chapter}{Conclusion}

L\textquotesingle installation d\textquotesingle iTop 3.2.3-1 sous Windows avec XAMPP a permis de mettre en place une plateforme locale complète de gestion des services informatiques. La configuration réalisée est adaptée à un environnement de test et d\textquotesingle apprentissage, tout en permettant d\textquotesingle expérimenter plusieurs processus ITIL ainsi que les fonctionnalités de la CMDB.


\chapter*{Perspectives}\addcontentsline{toc}{chapter}{Perspectives}

Pour prolonger ce travail, plusieurs pistes peuvent être envisagées, notamment l\textquotesingle installation et la configuration complète de Graphviz afin de bénéficier de l\textquotesingle ensemble des fonctionnalités graphiques d\textquotesingle iTop, ainsi que la migration de la plateforme vers un serveur Linux (Ubuntu, Debian ou Red Hat) pour un éventuel déploiement en environnement de production, offrant de meilleures performances, une sécurité renforcée et une maintenance facilitée.


\chapter*{Ouverture}\addcontentsline{toc}{chapter}{Ouverture}

En ouverture, ce travail pourrait être prolongé par une installation d\textquotesingle iTop directement sous Linux (pile LAMP), plus proche d\textquotesingle un environnement de production et permettant de résoudre nativement la limitation liée à Graphviz rencontrée dans ce rapport.

À plus long terme, la rédaction du rapport lui-même pourrait aussi être envisagée avec LaTeX, mieux adapté qu\textquotesingle un traitement de texte classique à la gestion des figures, des références et du suivi de version d\textquotesingle un document technique.

